\chapter{Controle de Vers{\~a}o com Git para Iniciantes (Opcional)}
\label{ap:git}

Conforme indicado no Cap{\'i}tulo~\ref{cap:desktop}, o controle de vers{\~a}o com Git {\'e} uma medida de seguran{\c c}a \textbf{recomendada, por{\'e}m facultativa}. Docentes de qualquer {\'a}rea podem utiliz{\'a}-lo sem conhecimentos pr{\'e}vios de programa{\c c}{\~a}o: o objetivo aqui n{\~a}o {\'e} colaborar em c{\'o}digo, mas simplesmente \textbf{guardar vers{\~o}es de arquivos} (planos de aula, provas, rubricas) e poder voltar no tempo se algo der errado.

\section{O Conceito: um Hist{\'o}rico com Fotografias}

Pense no Git como um sistema que tira uma ``fotografia'' da pasta da disciplina em determinados momentos, chamados \textbf{commits}. Cada fotografia {\'e} acompanhada de uma legenda (a mensagem de commit), por exemplo ``vers{\~a}o final da prova P1''. Se um arquivo for apagado ou estragado depois, {\'e} poss{\'i}vel restaurar qualquer fotografia anterior com um {\'u}nico comando.

\section{Passe 1: Instalar o Git}

\begin{itemize}
  \item \textbf{Windows}: Baixe e instale o instalador oficial em \url{https://git-scm.com/download/win}, mantendo as op{\c c}{\~o}es padr{\~a}o. Ao final, o menu Iniciar passa a oferecer o \textbf{Git Bash}, um terminal simples e amig{\'a}vel para os pr{\'o}ximos passos;
  \item \textbf{Linux}: Instale pelo gerenciador de pacotes da distribui{\c c}{\~a}o, por exemplo \texttt{sudo apt install git} (Ubuntu/Debian) ou \texttt{sudo dnf install git} (Fedora);
  \item \textbf{macOS}: Digite \texttt{git --version} no Terminal; se o Git n{\~a}o estiver presente, o pr{\'o}prio macOS oferece a instala{\c c}{\~a}o das ferramentas de linha de comando.
\end{itemize}

Para conferir se a instala{\c c}{\~a}o deu certo, abra o terminal e digite:
\begin{lstlisting}[language=bash, caption={Verificando a instala{\c c}{\~a}o do Git}]
git --version
\end{lstlisting}

\section{Passe 2: Configurar o Nome e o E-mail (uma {\'u}nica vez)}

O Git usa essas informa{\c c}{\~o}es para assinar cada fotografia. N{\~a}o {\'e} necess{\'a}rio criar conta em nenhum servi{\c c}o na internet; os dados ficam apenas no seu computador.

\begin{lstlisting}[language=bash, caption={Configura{\c c}{\~a}o inicial do Git (nome e e-mail)}]
git config --global user.name "Seu Nome"
git config --global user.email "seu.email@instituicao.edu.br"
\end{lstlisting}

\section{Passe 3: Tornar a Pasta da Disciplina um Reposit{\'o}rio}

Abra o terminal (\textbf{Windows}: Git Bash, clicando com o bot{\~a}o direito dentro da pasta da disciplina e escolhendo ``Git Bash Here''; \textbf{Linux/macOS}: terminal comum) e execute:

\begin{lstlisting}[language=bash, caption={Criando o reposit{\'o}rio da disciplina}]
git init
\end{lstlisting}

Esse comando cria, dentro da pasta, uma subpasta oculta \texttt{.git} que guardar{\'a} todo o hist{\'o}rico. N{\~a}o a apague.

\section{Passe 4: O Ciclo B{\'a}sico — Fotografar e Legendar}

A partir de agora, adote este ritual r{\'a}pido sempre que terminar uma vers{\~a}o de um material (ap{\'o}s criar ou editar a prova, o plano de aula, a rubrica):

\begin{lstlisting}[language=bash, caption={Salvando uma nova vers{\~a}o (commit)}]
git add .
git commit -m "versao final da prova P1"
\end{lstlisting}

\begin{itemize}
  \item \texttt{git add .} significa ``prepare todos os arquivos da pasta para a fotografia'';
  \item \texttt{git commit -m "..."} significa ``tire a fotografia e legende-a''.
\end{itemize}

\section{Passe 5: Recuperar uma Vers{\~a}o Anterior (se algo der errado)}

Para listar as fotografias j{\'a} tiradas, exibindo legenda, data e um c{\'o}digo identificador:
\begin{lstlisting}[language=bash, caption={Listando o hist{\'o}rico de vers{\~o}es}]
git log --oneline
\end{lstlisting}

Para restaurar um \'arquivo espec\'ifico ao estado da \'ultima fotografia (desfazendo edi{\c c}{\~o}es n{\~a}o desejadas):
\begin{lstlisting}[language=bash, caption={Restaurando um arquivo para a \'ultima vers{\~a}o salva}]
git restore prova-p1.md
\end{lstlisting}

Para voltar o \textbf{projeto inteiro} a uma fotografia anterior (por exemplo, a que aparece como \texttt{abc1234}):
\begin{lstlisting}[language=bash, caption={Restaurando o projeto a uma vers{\~a}o anterior}]
git checkout abc1234 -- .
\end{lstlisting}

\begin{dicaopencode}
\textbf{Seguran{\c c}a antes de tudo}: fa{\c c}a um \texttt{commit} antes de experimentos ousados (reformular inteiramente uma prova, remover uma se{\c c}{\~a}o do plano, etc.). Assim, qualquer resultado inesperado {\'e} revertido em um {\'u}nico comando, e a interface do OpenCode passa a exibir o hist{\'o}rico de vers{\~o}es diretamente no painel lateral, permitindo acompanhar visualmente o que mudou a cada passo.
\end{dicaopencode}
